\thispagestyle{abstractpagestyle}

\vspace*{36pt}

\begin{center}
    \textbf{\fontsize{20pt}{24pt} \selectfont ABSTRACT}
\end{center}

\vspace{24pt}

The proliferation of software-defined vehicle (SDV) platforms has opened an
opportunity to combine modern deep reinforcement learning (RL) with real
automotive runtime stacks. This thesis presents \textbf{SHILATE} --- a
complete, open-source pipeline that trains a Proximal Policy Optimization (PPO)
agent inside a Unity physics simulation and deploys the resulting policy onto an
Eclipse Leda SDV stack, bridging the gap between laboratory RL research and
realistic vehicle software deployment.

The simulation component is built in Unity 2022 LTS. A single rear-wheel-drive
vehicle navigates a circular obstacle course instrumented with a 21-ray
LiDAR-style sensor. The training bridge computes a shaped reward signal from
progress, collision events, and episode milestones, and communicates with the
Python training process exclusively over MQTT, keeping the two subsystems fully
decoupled. A dedicated Editor window (\textit{SHILATE Training Controller})
provides real-time health monitoring, metric graphs, and one-click training
control without leaving the Unity Editor.

The Python training stack wraps the simulation in a Gymnasium-compliant
environment, configures Stable-Baselines3 PPO with a two-layer actor-critic
network, and streams structured telemetry back to the Editor. Trained
checkpoints are subsequently deployed to the Eclipse Leda SDV image via a
containerised MQTT-to-Kuksa bridge and a Velocitas vehicle application that
enforces operational limits (overspeed, RPM redline) on the running agent.

Experimental results demonstrate that the PPO agent converges reliably on the
obstacle course within 100\,000 interaction steps and that the end-to-end
deployment chain correctly propagates vehicle telemetry from the simulation
through the VSS data layer to the Velocitas safety monitor. The project shows
that an integrated sim-to-SDV pipeline is achievable with commodity open-source
tooling, and identifies curriculum learning and multi-environment
parallelisation as natural extensions for future work.

\textbf{Keywords:} reinforcement learning, software-defined vehicle, Unity
simulation, PPO, MQTT, Eclipse Leda, Kuksa, Stable-Baselines3, Gymnasium.

\vfill
